\begin{figure}[]
  \centering
  \includegraphics[width=\linewidth]{Images/Concept Map.png}
  \caption{Concept Map}
  \label{fig:concept}
  \Description{}
\end{figure}


In the second week of our course, we embarked on our first group discussion. Our primary goal was to pinpoint a topic that intrigued everyone, and I suggested the theme of social connectivity, arguing that people place significant value on their social connections. Enhancing this aspect could offer users an enchanting experience, and even potentially leaving them with warm memories. My team members were all interested in this idea. We then proceeded to draft our concept map on the whiteboard, as shown in Figure \ref{fig:concept}. During our discussion, we outlined elements such as lighting, soft textures, and real-time tactile feedback to foster a sense of warmth. We categorized communication into verbal and non-verbal and decided to focus more on the latter to align with the "enchanted" theme. We also addressed the issue of time differences and the resulting misalignment in daily routines, a point deeply resonant with the three international students in our group. After completing the concept maps, we initiated our design synthesis process, culminating in a list of 10 ideas. This list helped us form a preliminary impression of what we aimed to create.

In the following lab session that week, we began refining the features our product should encompass. We specifically envisioned the "enchanted object" as a decorative item for the bedside, possibly a lamp or another object, due to its proximity to both the physical and emotional space of the user. Bastiaan initially suggested a snow globe, which I felt was not an everyday object (which proved wrong!), so we decided to further explore the product's functionalities. We unanimously agreed that continuously displaying the ambient light of the partner's device as a means to stay connected with their life was a compelling feature and could serve as a non-intrusive, companion-like device's core premise. Additionally, we considered incorporating the functionality of displaying the weather conditions of the partner's location. After establishing the product's tone, we shifted our focus to the "enchanted" aspect. I proposed that we should allow users to send a "signal" to each other, leading to discussions about what this signal could be. Our ideas included a touchpad and screen that would display the trajectories of the partner's touches in real-time, slowly fading over 3-4 hours. This would allow users to creatively use the touchpad and, over time, develop unique patterns symbolizing their bond. We also considered other signal options, such as sending songs to each other. These discussions formed our preliminary ideas. Although we had yet to finalize what the object would be, we gained a clearer understanding of the product we wanted to develop.